%% UCAS Homework starter -- homework.sty, adapted for the University of
%% Chinese Academy of Sciences from Josh Davis's well-known homework
%% template (the vendored .sty's own MIT header credits Josh Davis,
%% 2014).
%%
%% homework.sty IS BUNDLED here: it is not on CTAN, so there is nothing
%% for a system install to resolve.
%%
%% Source: https://github.com/HuangZhenyang/Latex-Template-for-UCAS-Homework
%% Licence: MIT, (c) 2014 Josh Davis (original template).
%%
%% engine.txt: xelatex (required -- the class \usepackage{fontspec} and
%% \usepackage[UTF8]{ctex} for CJK).
%%
%% THE FONT TRAP: the vendored homework.sty ORIGINALLY called
%% \setmonofont{Consolas}, a Windows-only commercial font absent from
%% this (and most Linux) TeX installations -- confirmed via both
%% fc-list and kpsewhich in the pinned image. Patched in homework.sty
%% to DejaVu Sans Mono, which resolves both ways and ships with every
%% modern TeX distribution. See the PATCH comment in homework.sty.
\documentclass{article}
\usepackage{homework}

%% This starter drops the upstream title page's UCAS logo image
%% (images/title/ucas-logo.png) so the download has no external asset
%% dependency -- plain text title instead.
\title{
    \textmd{\textbf{\hmwkClass\ \hmwkTitle}}\\
    \textmd{\textbf{\hmwkSubTitle}}\\
    \normalsize\vspace{0.1in}\small{\hmwkCompleteTime}\\
    \vspace{0.1in}\large{\textit{\hmwkClassInstructor\ }}\\
    \vspace{1in}
}

\renewcommand{\hmwkClass}{\textit{《算法设计与分析》}}
\renewcommand{\hmwkTitle}{Homework 4}
\renewcommand{\hmwkSubTitle}{Asymptotic Notation and Recurrences}
\renewcommand{\hmwkCompleteTime}{January 8, 2027}
\renewcommand{\hmwkClassInstructor}{Professor: Example Instructor}
\renewcommand{\hmwkAuthorName}{\textbf{Ada Lovelace} }
\renewcommand{\hmwkAuthorStuID}{\textbf{2027000001} }

\author{\hmwkAuthorName \\
	\hmwkAuthorStuID}
\date{}

\renewcommand{\part}[1]{\textbf{\large Part \Alph{partCounter}}\stepcounter{partCounter}\\}

\begin{document}

\maketitle

\begin{homeworkProblem}
	Give an appropriate positive constant $c$ such that $f(n) \leq c \cdot
	g(n)$ for all $n > 1$.

	\begin{enumerate}
		\item $f(n) = n^2 + n + 1$, $g(n) = 2n^3$
		\item $f(n) = n\sqrt{n} + n^2$, $g(n) = n^2$
		\item $f(n) = n^2 - n + 1$, $g(n) = n^2 / 2$
	\end{enumerate}

	\solution

	\begin{enumerate}
		\item For $n > 1$: $n^2 + n + 1 \le n^2 + n^2 + n^2 = 3n^2 < 2n^3$
		when $n > 1$, so $c = 1$ works for all $n > 1$.
		\item $n\sqrt{n} + n^2 \le n^2 + n^2 = 2n^2$ for $n \ge 1$, so $c=2$
		works.
		\item $n^2 - n + 1 \le n^2$ for $n \ge 1$, and $n^2 \le 2 \cdot
		(n^2/2)$, so $c = 2$ works.
	\end{enumerate}
\end{homeworkProblem}

\begin{homeworkProblem}
	Solve the recurrence $T(n) = 2T(n/2) + n$ using the recursion-tree
	method, and verify your answer using the Master Theorem.

	\solution

	At depth $i$ there are $2^i$ subproblems, each of size $n/2^i$, so
	the total work at depth $i$ is $2^i \cdot (n/2^i) = n$. The tree has
	$\log_2 n$ levels (down to subproblems of size 1), so
	\begin{equation}
		T(n) = \sum_{i=0}^{\log_2 n} n = n(\log_2 n + 1) = \Theta(n\log n).
	\end{equation}
	By the Master Theorem with $a = 2$, $b = 2$, $f(n) = n$: since
	$n^{\log_b a} = n^{\log_2 2} = n = f(n)$, this is Case 2, giving
	$T(n) = \Theta(n \log n)$, matching the recursion-tree result.
\end{homeworkProblem}

\begin{homeworkProblem}
	Sample pseudocode listing (demonstrates the monospace font used for
	code blocks):
	\begin{algorithm}
		\caption{MergeSort($A$, $p$, $r$)}
		\begin{algorithmic}[1]
			\If{$p < r$}
				\State $q \gets \lfloor (p+r)/2 \rfloor$
				\State MergeSort($A$, $p$, $q$)
				\State MergeSort($A$, $q+1$, $r$)
				\State Merge($A$, $p$, $q$, $r$)
			\EndIf
		\end{algorithmic}
	\end{algorithm}

	\solution

	The recurrence for merge sort is $T(n) = 2T(n/2) + \Theta(n)$, which
	by the same Master Theorem case as Problem 2 gives
	$T(n) = \Theta(n \log n)$.
\end{homeworkProblem}

\end{document}
